\section{Уроки 13--14: многодронные сценарии и итоговый проект}

\subsection{Цели уроков}
\begin{itemize}
\item Понять организацию роя через потоки и события.
\item Научиться синхронизировать этапы взлёта и переходов по точкам.
\item Спроектировать итоговую миссию с критериями приёмки.
\end{itemize}

\subsection{Опорные листинги}
\lstinputlisting[style=GeoskanStyle,language=pythoncustom,caption={Синхронный полёт двух дронов}]{code/python/examples/group_flight_examples/swarm_OPT.py}
\lstinputlisting[style=GeoskanStyle,language=pythoncustom,caption={Синхронный полёт с PiZero и MAVLink-портом}]{code/python/examples/group_flight_examples/swarm_OPT_PiZero.py}
\lstinputlisting[style=GeoskanStyle,language=pythoncustom,caption={Круговой роевой полёт с фазовым сдвигом}]{code/python/examples/group_flight_examples/swarm_circle_IRNAV.py}

\subsection{Итоговый проект}
\begin{infobox}[ТЗ проекта]
Команда разрабатывает миссию из трёх этапов: взлёт, фигура по точкам, посадка. На каждом этапе фиксируются телеметрия и один критерий безопасности.
\end{infobox}

\begin{table}[H]
\centering
\begin{tabularx}{\textwidth}{>{\raggedright\arraybackslash}p{0.26\textwidth}>{\raggedright\arraybackslash}X>{\raggedright\arraybackslash}p{0.23\textwidth}}
\toprule
Этап & Проверяемый результат & Критерий зачёта \\
\midrule
Старт миссии & Все аппараты перешли в состояние полёта & Нет рассинхронизации старта более 2 с \\
Маршрут & Прохождение всех точек & Не менее 90\% точек достигнуто \\
Завершение & Штатная посадка & Нет аварийных выключений \\
\bottomrule
\end{tabularx}
\caption{Рубрикатор итоговой практики}
\end{table}

\subsection{Контрольные вопросы}
\begin{enumerate}
\item Для чего в роевых скриптах используется \texttt{threading.Event}?
\item Почему в учебном коде нежелателен \texttt{while not point\_reached(): pass} без паузы?
\item Какие данные нужно логировать для разбора инцидентов?
\end{enumerate}
